\SourceLeaf{452}%
\begin{center}
\begin{tabular}{@{}l p{.54\linewidth} r@{}}
\(a.\)&\raggedright Length of 1505 measuring rods without correction \(.\ .\ .\ .\ .\ .\ .\ .\ .\)&\(3010{,}00000\) toises\\
\(b.\)&\raggedright Sum of the gaps measured with the glass wedges and the relevant full and half diameters of the plumbing cylinders \(.\)&\(+3{,}58389\) T.\\
\(c.\)&\raggedright Length of the supplementary rod \(.\ .\ .\ .\ .\ .\ .\ .\ .\ .\ .\ .\ .\ .\ .\)&\(+1{,}22106\) T.\\
\(d.\)&\raggedright Correction for the inclination of the plumbing cylinders to the plumb line \(.\)&\(-0{,}00008\) T.\\
\(e.\)&\raggedright Correction for the deviation of the rods from the alignment \(.\ .\ .\ .\ .\)&\(-0{,}00051\) T.\\
\(f.\)&\raggedright Correction for errors in the lengths of the measuring rods \(.\ .\ .\ .\ .\ .\ .\)&\(-0{,}10245\) T.\\
\(g.\)&\raggedright Correction for the temperature of the measuring rods differing from \(13^{0}\) R.&\(-0{,}19906\) T.\\
\(h.\)&\raggedright Reduction to the surface of the sea \(.\ .\ .\ .\ .\ .\ .\ .\ .\ .\ .\)&\(-0{,}02264\) T.\\\cline{3-3}
&\raggedright Length of the Braak baseline according to the new calculation \(.\ .\ .\ .\ .\)&\(=3014{,}48021\) toises
\end{tabular}
\end{center}

There is, however, a weak point in this calculation as well, namely the correction for the temperature of the measuring rods listed under \(g\). \textsc{Schumacher} published an illustrated description of the apparatus used to measure the baseline in ‘Schreiben an Dr. \textsc{Olbers} in Bremen etc. etc.’, Altona 1821’. It shows that the temperatures were measured with ordinary inset thermometers, rather than with metal thermometers. This is an unfortunate circumstance in itself, but the situation is much worse because the capacity of the rods to expand was inferred from only a few measurements of their lengths made in the field in the evening and in the morning. The difficulty can, however, be remedied. In 1853, rod No. IV. from \textsc{Schumacher}’s baseline apparatus was taken to Pulkowa for direct comparison with the standards collected there. The expansion of this rod for \(100^{0}\) was also obtained on that occasion. If one calculates the value of the coefficients of expansion given by \textsc{Struve} (see ‘Arc du méridien entre le Danube et la mer glaciale’, p. 51), for the correction under \(g\) one obtains \(-0{,}22812\) in place of \(-0{,}19906\).

With this adjustment, which Professor \textsc{Peters} also adopts, the length of the \textit{Braak baseline} is then
\[
=3014{,}451\ \text{toises},
\]
and this value must be regarded as \textit{definitive}. I should add that the expression of the length in these toises rests on comparison with the Pulkowa \textsc{Fortin} standard. Since that standard agrees with \textsc{Bessel}’s toise apart from a negligible difference, the length may also reasonably be regarded as expressed in \textsc{Bessel}’s toises.’

The coordinate register above gives a baseline length of \(5875{,}3614\) in the units used there, or \(3014{,}5757\) toises for an Earth meridian of \(360\times57009{,}746\) toises.

\textsc{Eduard Schmidt. Gauss} to \textsc{Schumacher}: Göttingen 1830 April 30. ‘To increase your confidence in \textsc{Schmidt}’s calculation, I should mention that he calculated the two principal elements of the dimensions of the Earth four times. — — He communicated result (IV) to me in manuscript. It is the same result that underlies my new auxiliary tables, namely flattening \(\frac{1}{297{,}732}\) and \(\frac{\text{Earth’s meridian quadrant}}{90}=57008^{\mathrm{T}}551\).’

\textsc{Bessel}. ‘Ueber einen Fehler in der Berechnung der französischen Gradmessung und seinen Einfluss auf die Bestimmung der Figur der Erde.’ Astronomische Nachrichten No. 438, vol. 19. p. 116. 1841 December 2. ‘Mean degree of the meridian \(=57013{,}109\) toises, semi-major axis \(a=3272077{,}14\) toises, semi-minor axis \(b=3261139{,}33\) toises, \(a:b=299{,}1528:298{,}1528\)’

When using the coordinates given in the registers above, they must therefore be multiplied for the time being by the following correction factor, until the baseline and the connecting triangles as far as \textit{Hamburg — Hohenhorn} have been measured anew:
